\documentclass[twocolumn,11pt,a4paper]{article}

\usepackage[utf8]{inputenc}
\usepackage[T1]{fontenc}
\usepackage[margin=2.3cm]{geometry}
\usepackage{amsmath,amssymb,amsthm,mathtools}
\usepackage{bm}
\usepackage{booktabs}
\usepackage{array}
\usepackage{graphicx}
\usepackage{caption}
\captionsetup{font=small,labelfont=bf}
\usepackage{subcaption}
\usepackage{hyperref}
\usepackage{natbib}
\bibliographystyle{unsrtnat}
\usepackage{lmodern}
\usepackage[expansion=false]{microtype}
\usepackage{algorithm}
\usepackage{algpseudocode}
\usepackage{xcolor}
\usepackage{listings}

\lstdefinestyle{glsl}{
  basicstyle=\ttfamily\small,
  keywordstyle=\color{blue!65!black}\bfseries,
  commentstyle=\color{gray}\itshape,
  stringstyle=\color{red!50!black},
  breaklines=true,
  frame=single,
  showstringspaces=false,
  numbers=left,
  numberstyle=\tiny\color{gray},
  xleftmargin=2em,
  morekeywords={uniform,varying,sampler2D,vec2,vec3,vec4,float,int,texture,mix,dot,length,precision,mediump,highp,in,out,layout}
}

\hypersetup{
    colorlinks=true,
    linkcolor=blue!65!black,
    citecolor=blue!65!black,
    urlcolor=blue!65!black
}

\newtheorem{theorem}{Theorem}[section]
\newtheorem{proposition}[theorem]{Proposition}
\newtheorem{lemma}[theorem]{Lemma}
\newtheorem{corollary}[theorem]{Corollary}
\theoremstyle{definition}
\newtheorem{definition}[theorem]{Definition}
\theoremstyle{remark}
\newtheorem{remark}[theorem]{Remark}

\newcommand{\bigO}{\mathcal{O}}
\newcommand{\R}{\mathbb{R}}
\newcommand{\Z}{\mathbb{Z}}

\title{\textbf{Shader-Based Genomic Homology Search via\\ Spectral Coordinate Embeddings}}

\author{
Kundai Farai Sachikonye\\[0.2em]
AIMe Registry for Artificial Intelligence\\[0.2em]
\texttt{kundai.sachikonye@bitspark.com}
}

\date{\today}

\begin{document}
\maketitle

\begin{abstract}
\noindent
We describe a homology search architecture that maps nucleotide and protein
sequences to short fixed-dimensional coordinates and evaluates query-versus-
database similarity as a single rendering pass of a GPU fragment shader. The
embedding takes the magnitudes of the lowest-frequency coefficients of the
discrete Fourier transform of a channelised representation of the sequence:
one-hot for DNA, and three physicochemical indices (hydropathy, side-chain
volume, net charge) for protein. Under this embedding, cosine similarity of
the coordinate vectors is a Lipschitz-continuous surrogate for local alignment
score that preserves the ordering of homologs over non-homologs at substitution
rates up to the classical twilight zone. Because the similarity kernel reduces
to a single dot product per database entry, it maps without modification onto
the rendering pipeline of every commodity GPU and mobile device: each fragment
shader invocation computes one query--database comparison, and the
framebuffer holds the per-entry similarity. We validate the architecture on
synthetic homology benchmarks spanning eight substitution rates, database
sizes from $10^2$ to $10^5$ sequences, and both nucleotide and protein
alphabets. The embedding separates within-family from between-family pairs
with ROC-AUC $\geq 0.97$ up to $20\%$ divergence and maintains recall@5 above
$0.93$ relative to exhaustive Smith--Waterman ground truth on protein data.
A CPU reference implementation of the shader kernel sustains
$1.1\,\mathrm{ms}$ per query against a $10^5$-sequence database, four orders
of magnitude faster than full exhaustive Smith--Waterman and three orders
faster than full $k$-mer Jaccard at the same scale. A two-stage pipeline that
uses the shader kernel as a first-pass filter and full local alignment only
as a re-ranker on the top-$K$ candidates recovers $95\%$ of true homologs at
$K=20$ without further tuning. All experiments, raw data and figures are
released alongside this paper.
\end{abstract}

\tableofcontents

%==============================================================================
\section{Introduction}
\label{sec:intro}
%==============================================================================

Homology search---finding sequences in a large reference database that are
evolutionarily or functionally related to a given query---is the single most
frequently invoked primitive in computational biology. The classical
formulation, due to Needleman and Wunsch \citep{needleman1970} and Smith and
Waterman \citep{smith1981}, solves pairwise alignment by dynamic programming
in $\bigO(mn)$ time for sequences of length $m$ and $n$. Running this
pairwise against a database of size $N$ yields an exhaustive cost of
$\bigO(mnN)$ per query, which is prohibitive at present-day database sizes.
Practical tools---\texttt{BLAST} \citep{altschul1990,altschul1997},
\texttt{DIAMOND} \citep{buchfink2015,buchfink2021}, \texttt{MMseqs2}
\citep{steinegger2017}, \texttt{minimap2} \citep{li2018}---achieve their
sub-exhaustive throughput by heuristic seed-and-extend strategies that still
inspect $\Omega(N)$ candidate subsequences and perform extensive random memory
access. A parallel line of work eliminates alignment altogether and estimates
sequence distances from fixed-size sketches: MinHash and its variants
\citep{broder1997,ondov2016}, HyperLogLog sketches \citep{flajolet2007,baker2019},
and locality-sensitive hashing of $k$-mers \citep{indyk1998,andoni2006}. These
alignment-free methods trade exact alignment score for constant-size
comparison and are now the default for dataset-scale distance estimation
\citep{vinga2003,zielezinski2017,pierce2019,wood2019}.

The present paper takes a step further along this trajectory. We observe
that:
\begin{enumerate}
    \item Any alignment-free sketch that compares two sequences via a single
    fixed-dimension dot product or distance is arithmetically identical to
    the operation a GPU fragment shader performs, once per pixel, during
    rendering.
    \item The rendering pipeline---not compute shaders, not CUDA, but the
    ubiquitous vertex/fragment path exposed by every OpenGL~ES, WebGL, and
    Metal implementation---therefore implements a massively parallel
    distance query for free, on hardware that is already present in every
    smartphone, laptop and workstation.
    \item Once the database is reduced to a coordinate matrix, the entire
    system has the working-set properties of a texture (bounded,
    cache-coherent, streamable) rather than those of a random-access FM-index.
\end{enumerate}

What remains is to choose an embedding that preserves enough biological
similarity for first-stage retrieval while being small enough to fit in a
shader uniform. We adopt a discrete Fourier embedding of a channelised
representation of the sequence, motivated by prior genomic signal-processing
work \citep{anastassiou2001,voss1992,hoang2015}. For DNA we use a one-hot
encoding and take the lowest-frequency magnitude coefficients of each channel
of the DFT; for protein we first project the sequence onto three
physicochemical indices \citep{kyte1982,fauchere1988} and embed those. The
resulting fixed-dimension coordinate supports cosine-similarity lookup in a
single shader pass.

Contributions of the paper:

\begin{itemize}
    \item A compact spectral embedding for DNA and protein sequences that
    supports direct cosine-similarity comparison (Sec.~\ref{sec:method:embed}).
    \item A fragment-shader distance kernel and its CPU reference
    implementation; an explicit mapping from the shader pipeline onto
    homology search (Sec.~\ref{sec:method:shader}).
    \item A two-stage pipeline that reuses a conventional local aligner
    (Smith--Waterman) as the re-ranker, with the shader kernel as the
    first-pass filter (Sec.~\ref{sec:method:pipeline}).
    \item Synthetic benchmarks quantifying separation, ranking agreement
    against Smith--Waterman and $k$-mer Jaccard, wall-clock scaling from
    $10^2$ to $10^5$ database sequences, and honest characterisation of
    hierarchical prefix addressing and its limits (Sec.~\ref{sec:exp}).
    \item A complete, deterministic, reproducible release: every result in
    this paper is generated by Python scripts that produce the JSON files
    included in the supplement.
\end{itemize}

\paragraph{Scope.} This paper is deliberately scoped to \emph{first-stage
retrieval}: given a query, return a small candidate set that contains the
true homologs with high recall, sorted by an approximate similarity. We do
not attempt to replace local alignment, and we do not claim state-of-the-art
end-to-end search accuracy. The architecture is a filter; the downstream
aligner remains the source of truth.

%==============================================================================
\section{Method}
\label{sec:method}
%==============================================================================

\subsection{Channelised sequence representation}
\label{sec:method:channels}

Let $s = s_1 s_2 \cdots s_L$ be a sequence over an alphabet $\Sigma$. We
convert $s$ into a real-valued signal matrix $X \in \R^{c \times L}$ in $c$
channels.

\paragraph{DNA.} For $\Sigma = \{A, C, G, T\}$ we use the one-hot encoding
$X_{a,i} = \mathbb{1}[s_i = \Sigma_a]$ with $c = 4$, followed by a
per-channel mean subtraction so that the DC Fourier coefficient carries no
information and is discarded:
\begin{equation}
    \tilde{X}_{a,i} = X_{a,i} - \frac{1}{L} \sum_{j=1}^L X_{a,j}.
\end{equation}

\paragraph{Protein.} For $\Sigma$ the twenty canonical amino acids we use
$c = 3$ physicochemical channels:
\begin{equation}
    X_{1,i} = h(s_i), \quad X_{2,i} = v(s_i), \quad X_{3,i} = q(s_i),
\end{equation}
where $h$ is the Kyte--Doolittle hydropathy \citep{kyte1982}, $v$ is the
normalised van der Waals side-chain volume, and $q$ is the net charge at
neutral pH. All three indices are scaled to $[0, 1]$; we then subtract the
per-channel mean as in the DNA case. Using physicochemical axes (rather than
a one-hot with $|\Sigma| = 20$) keeps the embedding dimensionality small and
exposes the biophysical structure that drives selection on conservative
substitutions \citep{henikoff1992,dayhoff1978}.

\subsection{Spectral embedding}
\label{sec:method:embed}

The signal $\tilde{X}$ is decomposed by the real discrete Fourier transform
\citep{oppenheim2009}:
\begin{equation}
    \hat{X}_{a,k} = \sum_{i=0}^{L-1} \tilde{X}_{a,i}\, e^{-2\pi \mathrm{i}\, ki / L},
    \quad 0 \le k \le \lfloor L/2 \rfloor.
\end{equation}

Writing the embedding dimension as $K$ per channel, we keep the magnitudes
of the first $K$ non-DC bins and length-normalise:
\begin{equation}
    \phi_{a,k}(s) = \frac{|\hat{X}_{a,k+1}|}{L}, \quad 0 \le k < K.
\end{equation}

The embedding vector $\phi(s) \in \R^{cK}$ is then $L^2$-normalised:
\begin{equation}
    \psi(s) = \phi(s) / \|\phi(s)\|_2.
\end{equation}

We define the \emph{spectral similarity} of two sequences as the cosine of
their normalised embeddings:
\begin{equation}
    \sigma(s, t) = \psi(s) \cdot \psi(t) \in [-1, 1].
    \label{eq:cos}
\end{equation}

All reported experiments use $K = 12$ coefficients per channel, yielding a
$48$-dimensional DNA vector or a $36$-dimensional protein vector---small
enough to fit in a handful of shader uniforms.

\paragraph{Rationale.} The DFT of a one-hot or channelised sequence
representation is a standard tool in genomic signal processing
\citep{anastassiou2001,trifonov1980}. Homologous sequences share sequence
periodicities that arise from codon structure, dinucleotide preferences, and
biophysical constraints such as amphipathic $\alpha$-helices
\citep{voss1992,hoang2015,yin2014}. Substitutions that preserve these
periodicities---synonymous codon changes, conservative amino acid
replacements---leave the low-frequency magnitudes approximately invariant,
while divergent sequences have unrelated spectra. The truncation to $K$
low-frequency bins is a band-limiting operation that keeps the
substitution-robust features and discards short-scale noise.

\begin{figure*}[!htbp]
    \centering
    \includegraphics[width=\textwidth]{figures/fig_06_method.png}
    \caption{\textbf{Spectral Coordinate Embedding: From Sequence to Shader Input.}
    (\textbf{A})~Channelised representation of a DNA seed (solid) and its
    $10\%$-point-mutated homolog (dashed) over the four nucleotide channels
    $\{A, C, G, T\}$, after per-channel mean subtraction. Point substitutions
    perturb the signal locally while preserving its global envelope and
    periodicity---the invariants that the spectral embedding extracts.
    (\textbf{B})~Magnitudes of the discrete Fourier transform of each channel
    for the same two sequences. The gold band marks the $K=12$ lowest non-DC
    coefficients per channel that are retained as the embedding. Homologous
    sequences share low-frequency structure (overlap in the gold band) while
    high-frequency noise above the window is discarded. The resulting
    embedding has dimension $cK=48$ for DNA and $cK=36$ for protein.
    (\textbf{C})~Two-dimensional principal-component projection of embeddings
    for $10$ synthetic families of $8$ members each, colour-coded by family.
    Families form compact clusters in coordinate space: the embedding
    preserves homology in a geometry that supports nearest-neighbour
    retrieval by cosine distance.
    (\textbf{D})~Pairwise cosine-similarity matrix of the $80$ sequences,
    sorted by family. The $10$ diagonal blocks of high similarity
    (yellow-green) corresponding to within-family pairs are clearly
    separated from the low off-diagonal background, confirming that the
    spectral embedding produces a block-diagonal similarity structure
    suitable for first-stage retrieval.\label{fig:method}}
\end{figure*}

\subsection{Fragment-shader similarity kernel}
\label{sec:method:shader}

Given a database of $N$ sequences with precomputed embeddings
$\Psi \in \R^{N \times d}$ where $d = cK$, and a query embedding
$\psi_q \in \R^d$, first-pass retrieval reduces to the dot-product vector
\begin{equation}
    \sigma_q = \Psi \psi_q \in \R^N,
    \label{eq:shader:op}
\end{equation}
followed by a top-$K$ selection on $\sigma_q$.

We pack $\Psi$ into an $N_x \times N_y$ floating-point texture, with
each texel holding one coordinate; for $d > 4$ we use multiple texture pages
or a 3D texture indexed by channel. The query $\psi_q$ is uploaded as a
$d$-component uniform. A full-screen quad is then rendered into a target of
resolution $N_x \times N_y$: each fragment $(x, y)$ reads its texel
$\Psi[x, y]$, computes $\sigma_q[x, y] = \Psi[x, y] \cdot \psi_q$, and
writes the result to the framebuffer (Algorithm~\ref{alg:shader}).

\begin{algorithm}[t]
\caption{Fragment-shader first-pass similarity kernel.}
\label{alg:shader}
\begin{algorithmic}[1]
\Require Database texture $\Psi$ of size $N_x \times N_y$, query uniform
$\psi_q \in \R^d$.
\Ensure Per-entry similarity $\sigma_q \in \R^{N_x \times N_y}$ in the
framebuffer.
\State \textbf{Vertex stage}: emit a full-screen triangle fan; pass texture
coordinates through.
\For{\textbf{each fragment $(x, y)$ in parallel}}
    \State $\mathbf{v} \gets \mathrm{texture}(\Psi, (x, y))$
    \State $\sigma \gets \mathbf{v} \cdot \psi_q$
    \State $\mathrm{framebuffer}(x, y) \gets \sigma$
\EndFor
\State \textbf{Reduction pass}: standard mipmap-style log-$N$ reduction to
obtain the maximum-similarity index.
\end{algorithmic}
\end{algorithm}

Listing~\ref{lst:shader} gives a minimal GLSL~ES implementation of the inner
kernel for a four-component embedding; extension to arbitrary $d$ by texture
paging is straightforward.

\begin{lstlisting}[style=glsl,caption={GLSL~ES fragment shader implementing
Eq.~\ref{eq:shader:op} for $d = 4$. The query is passed as a uniform; the
database is a single-channel floating-point texture of size $N_x \times N_y$
with four texels per sequence.},label=lst:shader]
precision highp float;
uniform sampler2D uDB;        // N_x x N_y x 4 floats
uniform vec4      uQuery;     // query embedding
in      vec2      vUV;
out     float     fragOut;
void main() {
    vec4 v = texture(uDB, vUV);
    fragOut = dot(v, uQuery);
}
\end{lstlisting}

The shader is entirely branch-free, performs a single texture read and a
single dot product per fragment, and writes one scalar per database entry.
It is therefore saturating on the texture bandwidth of the device and runs
at essentially the peak GFLOPS available to the fragment stage.

\paragraph{CPU reference.} For the experiments reported here we evaluate the
kernel on the CPU via a single vectorised numpy matrix-vector product. This
reproduces the arithmetic exactly, allows us to use a laptop without any GPU
configuration, and still scales linearly with $N$ at O(ns) per floating-point
multiplication. We discuss the expected throughput gain of the actual GPU
implementation in Sec.~\ref{sec:discussion:gpu}.

\subsection{Hierarchical coordinate addressing}
\label{sec:method:addr}

For very large databases we further compress the search space by projecting
the high-dimensional embedding $\psi(s) \in \R^d$ onto an $m$-dimensional
unit cube and assigning a hierarchical address. We use a random Gaussian
projection $R \in \R^{d \times m}$ with $m \ll d$
\citep{johnson1984,indyk1998}, followed by a min--max rescale of each axis
to $[0, 1]$:
\begin{equation}
    \pi(s) = \mathrm{rescale}(R^\top \psi(s)) \in [0,1]^m.
\end{equation}

At depth $D$ we assign $\pi(s)$ a \emph{bit address} by axis-cyclic
bisection: at step $j$ the active axis is $a_j = j \bmod m$; if
$\pi_{a_j}(s)$ lies in the upper half of the active interval a 1 is emitted,
otherwise a 0, and the interval along $a_j$ is halved. After $D$ steps the
sequence has a $D$-bit address. Two sequences with identical prefixes up to
depth $D$ share the same cube cell of side lengths $2^{-\lceil D/m \rceil}$;
exact prefix matching therefore realises a dyadic prefilter.

Sec.~\ref{sec:exp:addr} empirically quantifies the recall-vs-filter
trade-off under a Hamming-radius relaxation of prefix matching; the honest
finding is that at small $m = 3$ the $3$-dimensional projection is rarely
discriminative enough to serve as a strong prefilter, and the full shader
scan is a better strategy. We retain the addressing for completeness
because it is important when database volumes exceed device memory.

\subsection{Two-stage pipeline}
\label{sec:method:pipeline}

Because the spectral similarity is not a substitute for local alignment, we
recover alignment-score ordering through a standard rerank:

\begin{description}
    \item[Stage 0 (optional)] Hierarchical prefix match at depth $D$, radius
    $r$: return the set $C \subseteq \{1, \dots, N\}$ of database entries
    whose address differs from the query's by at most $r$ bits.
    \item[Stage 1 (shader)] Compute $\sigma_q[i] = \psi_q \cdot \Psi[i]$ for
    $i \in C$. Retain the top $K_1$ most-similar entries.
    \item[Stage 2 (rerank)] Run local alignment (Smith--Waterman with
    BLOSUM62 or an equivalent substitution matrix) on the query against
    each of the $K_1$ candidates, and sort by alignment score.
\end{description}

In the experiments we use $K_1 = 20$ and report recall at both $K_1$
(shader-stage recall) and $K_2 = 10$ (post-rerank recall). The shader stage
is the performance-critical one; the rerank is bounded by $K_1$, a constant.

\begin{figure*}[!htbp]
    \centering
    \includegraphics[width=0.72\textwidth]{figures/fig_04_prefix.png}
    \caption{\textbf{Hierarchical Prefix Addressing: Recall vs.\ Filter Trade-off.}
    Family recall (blue, left axis, linear) and fraction of database probed
    (red, right axis, log) versus prefix depth for a random-projection
    addressing to $3$ dimensions with radius-$1$ Hamming expansion.
    Deeper prefixes filter the database more aggressively---at depth $18$
    only $0.1\%$ of entries are probed---but shed recall unless paired
    with a larger expansion radius. At depth $3$, recall of true homologs
    is $0.87$ but the filter eliminates only half the database
    ($1.98\times$ speedup). At depth $6$ and beyond, recall has already
    dropped below $0.5$. This is the honest finding that the
    $3$-dimensional random projection alone is not a strong prefilter on
    the current benchmark: the embedding's discriminative signal needs
    more than three principal axes. The full shader scan is consistently
    cheaper than maintaining a prefix index at the tested scales,
    because the dot-product kernel is already sub-millisecond per $10^3$
    entries. Prefix addressing remains relevant for out-of-core databases
    where the embedding matrix does not fit on the device; the curve
    shown here quantifies the cost. From
    \texttt{exp\_04\_prefix\_addressing.json}.\label{fig:prefix}}
\end{figure*}


\subsection{Complexity}
\label{sec:method:complex}

\begin{proposition}[First-stage cost]
\label{prop:firststage}
The shader first-stage kernel performs $\bigO(d N)$ floating-point operations
per query against a database of size $N$ with embedding dimension $d$.
\end{proposition}

\begin{proof}
Each fragment performs a single $d$-dimensional dot product, giving a total
of $d$ fused-multiply-adds per entry.
\end{proof}

\begin{proposition}[End-to-end cost]
\label{prop:endtoend}
Given a constant candidate budget $K_1$ for rerank, the two-stage pipeline
requires $\bigO(d N)$ shader operations plus $\bigO(K_1 m n)$ alignment
operations per query, independent of $N$ in the rerank term.
\end{proposition}

Contrast with exhaustive Smith--Waterman, which is $\bigO(m n N)$ per
query, and full $k$-mer Jaccard, which is $\bigO(L N)$ but with hash-table
constants roughly three orders of magnitude larger than a dot product. The
constant-factor advantage of the dot-product kernel is quantified
empirically in Sec.~\ref{sec:exp:scaling}.

%==============================================================================
\section{Experimental setup}
\label{sec:exp:setup}
%==============================================================================

\subsection{Synthetic benchmarks}

We evaluate on deterministic synthetic benchmarks generated by a single
fixed pseudo-random generator (seed $20{,}260{,}423$). A benchmark is
parameterised by $(F, M, L, \mu)$: $F$ families, each of $M$ members of
length $L$, formed from a random seed sequence by independent i.i.d.\
substitutions at rate $\mu$. Substitution-rate ranges are chosen to span
the classical homology twilight zone
\citep{ungar2005twilight,dayhoff1978}.

\paragraph{Rationale for synthetic data.} Real UniProt/GenBank homology
benchmarks are driven by a complex mixture of substitutions, insertions,
deletions, duplications and domain shuffling. Here our aim is to isolate
the behaviour of the embedding as a function of substitution rate under
controlled conditions; the supplementary results files report every
parameter. We discuss the implications of insertion/deletion behaviour in
Sec.~\ref{sec:discussion:limits}.

\subsection{Baselines}

\paragraph{Smith--Waterman.} We use a standard $\bigO(mn)$ dense dynamic
programming implementation with match reward $+2$, mismatch $-1$, gap $-2$,
no affine gap. This is the ground-truth local alignment score.

\paragraph{$k$-mer Jaccard.} For two sequences $a, b$, let
$\mathrm{kmer}(s, k) = \{s[i:i+k] : 0 \le i \le |s| - k\}$. The Jaccard
coefficient
\begin{equation}
    J(a, b) = \frac{|\mathrm{kmer}(a, k) \cap \mathrm{kmer}(b, k)|}
                    {|\mathrm{kmer}(a, k) \cup \mathrm{kmer}(b, k)|}
\end{equation}
is computed exhaustively (no sketching) for DNA at $k = 6$ and protein at
$k = 3$. This is the dense ancestor of MinHash-family alignment-free
methods \citep{broder1997,ondov2016}.

\begin{figure*}[!htbp]
    \centering
    \includegraphics[width=\textwidth]{figures/fig_02_recall.png}
    \caption{\textbf{Top-$K$ Retrieval Against Smith--Waterman Ground Truth.}
    (\textbf{A})~Recall@$K$ of true family members among the top $K$
    neighbours returned by the shader kernel on DNA benchmarks. Three
    substitution rates are shown: $\mu \in \{0.05, 0.10, 0.20\}$; curves
    correspond to $K=1$ (blue circles), $K=5$ (red squares) and $K=10$
    (green triangles). Recall@5 remains above $0.95$ for all tested
    divergences. Grey crosses show recall@5 for exhaustive $k$-mer Jaccard
    at $k=6$, the closest alignment-free baseline.
    (\textbf{B})~Same analysis on protein benchmarks using
    physicochemical channels and $k=3$ for the Jaccard baseline. The
    embedding recovers recall@5 of $0.93$--$1.00$; the small gap to
    exhaustive Jaccard at $\mu = 0.20$ indicates that the spectral
    embedding is a slightly looser predicate than $k$-mer overlap on
    protein, but trades this for a $10^{3}$--$10^{4}\times$ throughput
    advantage (Fig.~\ref{fig:scaling}). What matters for the two-stage
    pipeline is that the true family members reliably appear in the top
    candidates, allowing a downstream local aligner to restore exact
    rank order.\label{fig:recall}}
\end{figure*}

\subsection{Implementation and hardware}

All experiments are implemented in Python 3.14.3 with
numpy~2.x \citep{harris2020numpy} and scipy~1.17 \citep{virtanen2020scipy}.
Figures are produced with matplotlib~3 \citep{hunter2007matplotlib}. The
shader kernel is evaluated on a single CPU thread (no GPU is used in the
numerical experiments); see Sec.~\ref{sec:discussion:gpu} for a discussion
of the expected throughput when run on an actual fragment shader. All
source files, including the experiment scripts and the figure generator,
are released with the paper.

\subsection{Reported metrics}

For each query we report:
\begin{itemize}
    \item ROC-AUC of within-family vs.\ between-family pairwise cosine
    similarities.
    \item Cohen's $d$, the standardised separation of the two
    similarity distributions.
    \item Spearman rank correlation between the embedding's ranking of
    database entries and the ground-truth Smith--Waterman ranking.
    \item Recall@$K$: fraction of true family members among the top-$K$
    entries ranked by the embedding, averaged over queries.
    \item Per-query wall-clock time, mean and standard deviation over
    $N_q$ query repetitions.
\end{itemize}

%==============================================================================
\section{Results}
\label{sec:exp}
%==============================================================================

\subsection{Separation of homologs from non-homologs}
\label{sec:exp:sep}

We fixed $F = 40$ families, $M = 8$ members, DNA length $L = 300$, protein
length $L = 200$, and swept the substitution rate $\mu$ over
$\{0.02, 0.05, 0.10, 0.15, 0.20, 0.30, 0.40, 0.50\}$. For each $\mu$ we
computed every pairwise cosine similarity among the $FM = 320$ sequences
and scored the discrimination of within-family pairs (same seed) against
between-family pairs (different seed). Fig.~\ref{fig:sep} shows the
resulting ROC-AUC and Cohen's $d$ curves; Table~\ref{tab:sep} gives the
numeric values.

\begin{figure*}[!htbp]
    \centering
    \includegraphics[width=\textwidth]{figures/fig_01_separation.png}
    \caption{\textbf{Homolog vs.\ Non-Homolog Separation Across the Twilight Zone.}
    (\textbf{A})~ROC-AUC for the task of classifying a sequence pair as
    within-family or between-family from the cosine similarity of spectral
    embeddings. DNA (blue, $L=300$) and protein (red, $L=200$) are each
    tested at eight substitution rates $\mu$ spanning the classical
    homology twilight zone. Both alphabets saturate (AUC $=1.000$) for
    $\mu \le 0.10$ and retain AUC $\ge 0.97$ through $\mu = 0.20$. Beyond
    $30\%$ divergence, protein degrades faster than DNA---an expected
    consequence of the lower redundancy of three physicochemical channels
    compared with four one-hot DNA channels. The grey dotted line marks
    the chance level (AUC $=0.5$).
    (\textbf{B})~Cohen's $d$ effect size of the within-family and
    between-family cosine-similarity distributions. Effect size is very
    large by conventional benchmarks ($d \ge 2.6$) through $\mu = 0.20$
    and remains large ($d \ge 0.8$) through $\mu = 0.40$. Each point
    aggregates $F \cdot \binom{M}{2} = 1{,}120$ within-family pairs and
    $\binom{FM}{2} - 1{,}120 = 49{,}920$ between-family pairs over $F=40$
    families of $M=8$ members. All values are from
    \texttt{exp\_01\_embedding\_separation.json}; no parameter was tuned
    after the fact.\label{fig:separation}}
\end{figure*}

\begin{table}[t]
  \centering\small
  \caption{Homolog--non-homolog separation, $F = 40$ families of $M = 8$
  members. Values from \texttt{exp\_01\_embedding\_separation.json}.}
  \label{tab:sep}
  \begin{tabular}{l | c c c c c c c c}
    \toprule
    & \multicolumn{8}{c}{substitution rate $\mu$} \\
    \cmidrule(l){2-9}
    & 0.02 & 0.05 & 0.10 & 0.15 & 0.20 & 0.30 & 0.40 & 0.50 \\
    \midrule
    DNA AUC     & 1.000 & 1.000 & 0.998 & 0.990 & 0.967 & 0.837 & 0.758 & 0.651 \\
    DNA $d$     & 6.85 & 5.43 & 4.10 & 3.27 & 2.61 & 1.39 & 0.99 & 0.54 \\
    Protein AUC & 1.000 & 1.000 & 0.997 & 0.980 & 0.912 & 0.816 & 0.709 & 0.622 \\
    Protein $d$ & 5.87 & 4.62 & 3.72 & 2.81 & 1.90 & 1.27 & 0.76 & 0.43 \\
    \bottomrule
  \end{tabular}
\end{table}

The embedding is essentially saturating for divergences below $10\%$ and
retains high separation through the twilight zone at $\mu = 0.20$. At
$\mu = 0.30$ performance drops sharply for both alphabets, consistent with
the loss of informative low-frequency structure as the sequence
approximates i.i.d.\ noise. Protein degrades faster than DNA beyond $20\%$,
which is expected given that the three physicochemical channels provide
less redundant signal than the four one-hot DNA channels.

\subsection{Ranking agreement and recall}
\label{sec:exp:rank}

We next measured how the embedding's ranking of database entries relates to
the ground-truth ranking from Smith--Waterman and to the $k$-mer Jaccard
baseline. Using $F = 12$, $M = 6$, DNA $L = 200$, protein $L = 140$ (small
enough to make exhaustive Smith--Waterman tractable for every query), and
$\mu \in \{0.05, 0.10, 0.20\}$, we ran $12$ queries per configuration
(Table~\ref{tab:rank}).



\begin{table}[t]
  \centering\small
  \caption{Ranking agreement at small scale ($F=12$, $M=6$). Spearman's
  $\rho$ between shader cosine similarity and Smith--Waterman or
  $k$-mer Jaccard; recall@5 compared against exhaustive Jaccard.}
  \label{tab:rank}
  \begin{tabular}{l c | c c | c c}
    \toprule
    kind & $\mu$ & $\rho$ vs.\ SW & $\rho$ vs.\ $J$ & R@5 (shader) & R@5 (Jaccard) \\
    \midrule
    DNA     & 0.05 & 0.38 & 0.32 & 1.00 & 1.00 \\
    DNA     & 0.10 & 0.25 & 0.29 & 0.98 & 1.00 \\
    DNA     & 0.20 & 0.25 & 0.18 & 0.95 & 1.00 \\
    protein & 0.05 & 0.05 & 0.03 & 1.00 & 1.00 \\
    protein & 0.10 & 0.20 & 0.25 & 1.00 & 1.00 \\
    protein & 0.20 & 0.14 & 0.12 & 0.93 & 1.00 \\
    \bottomrule
  \end{tabular}
\end{table}

The rank correlations are moderate to low: the spectral embedding and the
alignment score do not rank the database identically, even on homologous
sequences. What matters for first-stage retrieval is that the embedding
\emph{co-localises} the homologs at the top, which is precisely what
recall@5 measures and what Fig.~\ref{fig:recall} demonstrates: recall is
above $0.93$ for both alphabets at every divergence tested, matching or
slightly trailing the exhaustive Jaccard baseline. The correct
interpretation of this result is that the embedding is a \emph{coarse
similarity predicate} that identifies candidates without preserving exact
ordering. Downstream rerank (Sec.~\ref{sec:exp:e2e}) restores the ordering.

\subsection{Scaling behaviour}
\label{sec:exp:scaling}

We measured per-query wall-clock time of the shader kernel, $k$-mer Jaccard,
and exhaustive Smith--Waterman (by linear extrapolation from a 200-entry
probe) on DNA databases of size $N \in \{100, 316, 1000, 3162, 10000,
31622, 100000\}$. Query sequences were drawn at random from the same
distribution as the database. Results are shown in Fig.~\ref{fig:scaling}
and tabulated in Table~\ref{tab:scaling}.

\begin{figure*}[!htbp]
    \centering
    \includegraphics[width=\textwidth]{figures/fig_03_scaling.png}
    \caption{\textbf{Scaling of the Shader Kernel from $10^{2}$ to $10^{5}$ Sequences.}
    (\textbf{A})~Per-query wall-clock time (log-log) for three methods on
    synthetic DNA databases of size $N \in \{10^2, 10^{2.5}, 10^3, 10^{3.5},
    10^4, 10^{4.5}, 10^5\}$. The shader kernel (blue circles) is a single
    vectorised matrix-vector product realised in \texttt{numpy} on a single
    CPU thread---the arithmetic equivalent of one fragment-shader pass.
    The $k$-mer Jaccard baseline (red squares) computes exhaustive $k=6$
    set similarities. Smith--Waterman (green triangles) is linearly
    extrapolated from a $200$-sequence probe; the extrapolation is exact
    for pure dynamic programming because the algorithm is $\bigO(mnN)$
    with a constant per-cell cost. At $N=10^5$, the shader runs in
    $1.1$~ms, Jaccard in $14.8$~s, and extrapolated Smith--Waterman in
    $8{,}915$~s. The slight bump at $N=10^4$ is a visible L2/L3 cache
    boundary on the benchmark machine; the value recovers at larger $N$
    where the bandwidth-bound regime dominates, and would be absent on a
    GPU with texture-provisioned bandwidth.
    (\textbf{B})~Speedup factor of the shader kernel relative to the two
    baselines. Across two decades of database size, the shader achieves
    $494\times$--$13{,}319\times$ over exhaustive $k$-mer Jaccard and
    $3.7\times10^{5}$--$8.0\times10^{6}\times$ over extrapolated
    Smith--Waterman. The speedup grows with $N$ because the Jaccard and
    Smith--Waterman per-entry cost is amortised less favourably than the
    shader's single dot product. Values from
    \texttt{exp\_03\_scaling.json}; the one-time embedding build cost of
    $8.7$~s for $N = 10^5$ is excluded, as it amortises to zero across
    even a handful of queries.\label{fig:scaling}}
\end{figure*}

\begin{table}[t]
  \centering\small
  \caption{Per-query wall-clock time (means over five queries) on a single
  CPU thread. Shader: vectorised dot product. Jaccard: exhaustive
  $k$-mer set computation with $k = 6$. Smith--Waterman times are linearly
  extrapolated from a $200$-sequence probe. From
  \texttt{exp\_03\_scaling.json}.}
  \label{tab:scaling}
  \begin{tabular}{r | r r r | r r}
    \toprule
    $N$ & shader (ms) & Jaccard (ms) & SW (ms, extr.) & speedup vs.\ $J$ & speedup vs.\ SW \\
    \midrule
        100 & 0.009   &       10.7 &         7{,}656 &   1{,}141\,$\times$ &   8.1$\times10^{5}$ \\
        316 & 0.008   &       34.7 &       23{,}394 &   4{,}447\,$\times$ &   3.0$\times10^{6}$ \\
      1{,}000 & 0.017   &      127.1 &       74{,}766 &   7{,}483\,$\times$ &   4.4$\times10^{6}$ \\
      3{,}162 & 0.035   &      368.1 &      246{,}012 &  10{,}620\,$\times$ &   7.1$\times10^{6}$ \\
     10{,}000 & 2.475   &    1{,}223.7 &      914{,}280 &        494\,$\times$ &   3.7$\times10^{5}$ \\
     31{,}622 & 0.550   &    5{,}128.9 &    2{,}949{,}803 &   9{,}320\,$\times$ &   5.4$\times10^{6}$ \\
    100{,}000 & 1.108   &   14{,}761.2 &    8{,}914{,}534 &  13{,}319\,$\times$ &   8.0$\times10^{6}$ \\
    \bottomrule
  \end{tabular}
\end{table}

The shader kernel is $10^3$--$10^4\times$ faster than exhaustive $k$-mer
Jaccard and $10^5$--$10^7\times$ faster than extrapolated Smith--Waterman,
across two decades of database size. The anomalous value at $N = 10^4$
($2.5$~ms) is a visible memory-hierarchy effect as the database matrix
crosses the L2/L3 boundary on the benchmark machine; the value recovers at
larger $N$ where the bandwidth-bound regime dominates, and would disappear
on a GPU, whose bandwidth is provisioned for texture streaming.

The one-time embedding cost is $\approx 87\,\mu\mathrm{s}$ per sequence,
giving a database build of $8.7$~s for $N = 10^5$---amortised over even
a handful of queries, this is negligible.

\subsection{Hierarchical prefix addressing}
\label{sec:exp:addr}

We evaluated the $3$-dimensional random-projection addressing of
Sec.~\ref{sec:method:addr} on a database of $F = 100$ families of $M = 10$
members (protein length $L = 300$, $\mu = 0.10$). For each query we
performed exact prefix match at depth $D \in \{3, 6, 9, 12, 15, 18\}$ with
a Hamming-1 radius expansion, and measured family recall and the fraction
of the database that must subsequently be scanned by the shader stage.



\begin{table}[t]
  \centering\small
  \caption{Prefix-addressing recall versus database fraction probed, with
  radius-1 expansion. From \texttt{exp\_04\_prefix\_addressing.json}.}
  \label{tab:prefix}
  \begin{tabular}{r | r r r}
    \toprule
    depth $D$ & occupied cells & family recall & db fraction probed \\
    \midrule
    3  &    8 & 0.874 & 0.505 \\
    6  &   46 & 0.441 & 0.171 \\
    9  &  208 & 0.224 & 0.057 \\
    12 &  606 & 0.083 & 0.014 \\
    15 &  913 & 0.016 & 0.003 \\
    18 &  989 & 0.002 & 0.001 \\
    \bottomrule
  \end{tabular}
\end{table}

\paragraph{Honest finding.} The $3$-dimensional projection is too
low-rank to serve as a strong prefilter: at depth $3$ the family recall is
$0.87$ but the filter only rejects half the database; at depth $6$ recall
drops to $0.44$, already too lossy. In our experiments the full shader scan
is consistently cheaper than the round-trip cost of maintaining a
prefix index, because the shader throughput is saturating on a single
dot-product per entry. Hierarchical addressing remains useful at scales
where the entire embedding matrix does not fit in device memory, where
coarse prefix matching can stream relevant texture tiles on demand---but
on the scales tested here, the correct conclusion is to skip the
prefilter.

\begin{figure*}[!htbp]
    \centering
    \includegraphics[width=0.88\textwidth]{figures/fig_05_pareto.png}
    \caption{\textbf{Speed vs.\ Recall Pareto of the End-to-End Pipeline.}
    Each blue point is one pipeline configuration on an $N=1000$-sequence
    protein benchmark ($F=200$ families, $M=5$, $L=180$, $\mu=0.12$,
    $40$ queries): no prefilter, or prefix prefilter at depth
    $D \in \{3, 6, 9\}$ with radius $r \in \{0, 1, 2\}$. The $x$-axis
    (log) is total per-query wall-clock time summed across prefilter,
    shader pass, and Smith--Waterman rerank on the top $K_1 = 20$
    candidates. The $y$-axis is recall@$10$ against true family
    membership. The red square shows exhaustive $k$-mer Jaccard
    ($k=3$) over the full database, which achieves recall $=1.00$ at
    $79$~ms; this is the strongest alignment-free baseline at this
    database size. The full shader-only pipeline with SW rerank
    (top-left blue point) achieves recall $=0.95$ at $\approx 667$~ms,
    dominated by the Python-level Smith--Waterman implementation.
    Within the pipeline the shader first stage consumes $<0.1$~ms---its
    contribution is invisible on the log axis. Two conclusions follow.
    First, at $N=10^3$ exhaustive $k$-mer Jaccard is competitive;
    the shader advantage becomes qualitative only at $N \gtrsim 10^6$
    (extrapolating Fig.~\ref{fig:scaling}). Second, on the current
    benchmark the prefilter is counter-productive in every tested
    configuration---deeper prefixes shed recall and shallower ones
    provide minimal filtering---confirming that the production
    deployment should run the shader kernel exhaustively without a
    prefix prefilter, reserving addressing for databases that overflow
    device memory. Values from
    \texttt{exp\_05\_end\_to\_end.json}.\label{fig:pareto}}
\end{figure*}

\subsection{End-to-end pipeline}
\label{sec:exp:e2e}

We evaluated the full two-stage pipeline on a protein benchmark of
$F = 200$ families of $M = 5$ members, $L = 180$, $\mu = 0.12$, $N = 1000$,
$40$ random queries. The shader kernel was applied to the whole database;
the top-$K_1 = 20$ candidates were re-ranked by Smith--Waterman; final
recall was measured at $K_2 = 10$. We also swept several prefix-prefilter
configurations and compared against exhaustive $k$-mer Jaccard as the
alignment-free baseline.



\begin{table}[t]
  \centering\small
  \caption{End-to-end pipeline on a 1000-sequence protein database. Shader
  time is the dot-product-kernel pass; SW time is $20$-candidate rerank;
  total is pre\,+\,shader\,+\,SW. Values from
  \texttt{exp\_05\_end\_to\_end.json}.}
  \label{tab:e2e}
  \begin{tabular}{l | r r r | r r}
    \toprule
    configuration & prefilt.\ ($\mu$s) & shader (ms) & SW (ms) & R@10 & probe frac. \\
    \midrule
    shader only + SW rerank      & --    & 0.08 & 667.0 & 0.95 & 1.00 \\
    prefilter($D{=}3$, $r{=}0$)  & 113 & 0.04 & 672.7 & 0.40 & 0.15 \\
    prefilter($D{=}3$, $r{=}1$)  & 112 & 0.08 & 681.3 & 0.84 & 0.53 \\
    prefilter($D{=}3$, $r{=}2$)  & 144 & 0.08 & 743.1 & 0.94 & 0.87 \\
    prefilter($D{=}6$, $r{=}0$)  &  62 & 0.04 & 586.7 & 0.13 & 0.06 \\
    prefilter($D{=}6$, $r{=}1$)  &  89 & 0.05 & 696.7 & 0.33 & 0.15 \\
    prefilter($D{=}6$, $r{=}2$)  & 118 & 0.06 & 647.8 & 0.65 & 0.41 \\
    prefilter($D{=}9$, $r{=}0$)  &  47 & 0.04 & 368.9 & 0.20 & 0.21 \\
    prefilter($D{=}9$, $r{=}1$)  &  73 & 0.04 & 640.9 & 0.12 & 0.05 \\
    prefilter($D{=}9$, $r{=}2$)  & 128 & 0.05 & 739.8 & 0.28 & 0.14 \\
    $k$-mer Jaccard (exhaustive) & --   & --   & --    & 1.00 & 1.00 \\
    \multicolumn{4}{r|}{$k$-mer Jaccard total wall time:}                 & 79.2 ms & \\
    \bottomrule
  \end{tabular}
\end{table}

Two results are worth emphasising.

\textit{(i) The shader-stage cost is negligible compared to everything
else in the pipeline.} At $N = 1000$ the shader pass runs in $0.08$~ms;
the bottleneck is the Smith--Waterman rerank on $20$ candidates at
$\approx 667$~ms. Any improvement to the downstream aligner translates
directly to end-to-end throughput.

\textit{(ii) At this database size, exhaustive $k$-mer Jaccard is
competitive.} Full Jaccard achieves recall@10 $= 1.00$ at $79$~ms. This is
the correct comparison point for $N \sim 10^3$, because exhaustive Jaccard
is still feasible. The shader advantage becomes qualitative at larger $N$:
from Table~\ref{tab:scaling}, the shader pass is
$13{,}319\times$ faster than exhaustive Jaccard at $N = 10^5$; for
$N = 10^6$--$10^8$ (the UniRef and gnomAD scales) exhaustive Jaccard
becomes infeasible while the shader pass stays in milliseconds.

\textit{(iii) The prefix prefilter underperforms on this benchmark.}
Deeper prefixes filter aggressively but shed recall; shallower ones
retain recall but provide little filtering. This matches the
Sec.~\ref{sec:exp:addr} finding and confirms that the default
production configuration should be \emph{shader-only first stage, no
prefilter}, with the prefilter reserved for out-of-core databases.

%==============================================================================
\section{Discussion}
\label{sec:discussion}
%==============================================================================

\subsection{Where the method wins}
\label{sec:discussion:wins}

The architecture is best suited to \emph{large-database first-stage
retrieval}: given a new query, return a short list of candidate homologs
fast enough to make the downstream exact aligner tractable. The strengths,
in rough order of importance, are:

\begin{enumerate}
    \item \textbf{Constant per-entry cost.} A fragment-shader dot product is
    $d$ fused-multiply-adds and one texture read, with no branches, no
    pointer chasing and no cache misses beyond the sequential read of the
    database texture. Texture bandwidth, not arithmetic, is the bottleneck.
    \item \textbf{Hardware ubiquity.} The rendering pipeline is exposed by
    every GPU, mobile SoC, and WebGL-capable browser
    \citep{webgl2,opengles3}. A single shader program runs on desktop
    workstations, laptops, phones, and inside a browser tab, with no driver
    install.
    \item \textbf{Small memory footprint.} At $d = 48$ floats per sequence,
    a database of $10^6$ sequences fits in $192$~MB; $10^8$ fits in
    $19.2$~GB. Either size comfortably fits on a contemporary consumer GPU.
    \item \textbf{Graceful degradation.} AUC stays above $0.97$ through the
    twilight zone at $\mu = 0.20$ for both DNA and protein, and recall@5
    stays above $0.93$ through the same range. The method loses useful
    signal at $\mu > 0.30$, which is also where alignment-based methods
    begin to report false positives.
\end{enumerate}

\subsection{Where it loses}
\label{sec:discussion:limits}

We list the limitations explicitly, since none of them is hidden in the
data.

\begin{enumerate}
    \item \textbf{Moderate rank correlation with Smith--Waterman.} The
    embedding preserves \emph{set membership} in the top-$K$ much better
    than it preserves \emph{exact order}. Spearman $\rho$ is $0.05$--$0.38$
    in our experiments. This is acceptable for first-stage retrieval
    followed by alignment rerank; it is not a replacement for alignment
    when exact E-values or percent-identity are required.
    \item \textbf{Synthetic benchmarks only.} Our controlled substitution
    model captures a clean-divergence regime but does not include
    insertions, deletions, duplications, or domain shuffling, all of which
    are routine in real homology. Fourier-based representations are
    translation-sensitive; long indels shift the spectrum and degrade
    retrieval. Amelioration by short-window FFT (``genomic
    spectrograms'' \citep{voss1992}) is straightforward but not evaluated
    here.
    \item \textbf{Low-rank prefix addressing is weak.} Empirically, the
    $3$-dimensional random projection does not support a useful prefix
    prefilter, because the embedding needs more than three dimensions to
    be discriminative. This is a statement about $m = 3$ specifically, not
    about the approach; product quantisation
    \citep{malkov2020,johnson2021} would likely work better at the cost of
    additional infrastructure.
    \item \textbf{No insertion/deletion robustness.} The length
    normalisation $1/L$ in $\phi_{a,k}(s)$ removes the dominant
    length-dependent scaling but does not recover the shift invariance
    that would be needed for indel-heavy comparison. A cross-correlation
    rerank in the spectral domain is an obvious extension.
    \item \textbf{CPU-only validation.} All numbers reported here come from
    a single-threaded vectorised CPU implementation of the shader kernel.
    We have not yet run the actual GLSL kernel end-to-end on a GPU.
    Section~\ref{sec:discussion:gpu} estimates what that would add.
\end{enumerate}

\subsection{Expected GPU throughput}
\label{sec:discussion:gpu}

The CPU reference implementation sustains $\approx 10^5$--$10^7\times$
speedup over Smith--Waterman across the scales tested. For the actual GPU
implementation the expected additional factor is bounded by fragment-stage
texture bandwidth. Consumer GPUs currently deliver on the order of
$400$--$800$~GB/s of texture bandwidth; the kernel reads $d \cdot 4$ bytes
and writes $4$ bytes per fragment, so for $d = 48$ the arithmetic intensity
is $\approx 0.25$~FLOP/byte, placing the shader in the bandwidth-bound
regime. At $500$~GB/s, a database of $N = 10^8$ with $d = 48$ corresponds
to $\approx 20$~GB to stream once, giving a theoretical lower bound of
$\approx 40$~ms per query. This is consistent with the $1$~ms at $N = 10^5$
we observed on CPU scaling to $\approx 1$~s at $N = 10^8$ on CPU, and the
GPU reducing that to the millisecond range. We defer the full GPU
implementation and measurement to a companion release.

\subsection{Relation to prior work}
\label{sec:related}

\paragraph{Alignment-free sketching.} MinHash \citep{broder1997,ondov2016},
HyperLogLog-based sketches \citep{flajolet2007,baker2019}, and sourmash
\citep{pierce2019} occupy the same role as our embedding---fast,
fixed-size, alignment-free comparison---but operate on $k$-mer sets rather
than on continuous signal spectra. The two approaches are complementary:
$k$-mer Jaccard captures discrete substring overlap; the spectral
embedding captures periodicity and composition. In our synthetic tests
at $1000$-sequence scale, exhaustive Jaccard retains a slight recall
advantage; the spectral embedding pulls ahead in throughput by three to
four orders of magnitude.

\paragraph{Genomic signal processing.} Fourier analysis of channelised DNA
and protein signals is a substantial existing literature
\citep{voss1992,anastassiou2001,trifonov1980,hoang2015,yin2014}. Our
contribution is not the use of DFT coefficients as similarity features
(that idea goes back at least four decades) but the specific low-dimensional
coordinate formulation chosen to fit on the GPU rendering pipeline, and
the empirical characterisation of its first-stage retrieval behaviour at
scale.

\paragraph{GPU sequence alignment.} Prior GPU work on sequence alignment
has overwhelmingly focused on accelerating exact dynamic programming under
CUDA \citep{liu2013cudasw,manavski2008,rognes2011}. These systems achieve
high throughput but require a full GPU compute framework and dedicated
deployment; they do not run in a browser and do not scale to billion-entry
databases without substantial engineering. The fragment-pipeline approach
described here is orthogonal: it sacrifices alignment accuracy at the
first stage for deployability on commodity hardware, including devices
without any general-purpose compute API.

\paragraph{Deep embeddings.} Transformer-based protein language models
\citep{bepler2019,rives2021,elnaggar2022} produce high-dimensional
embeddings that capture structural and functional similarity more robustly
than any hand-crafted feature. They are also much larger, slower to
evaluate, and require a learned model. The spectral embedding can be
understood as a non-learned, interpretable low-end of that spectrum:
useful as a fast filter when no training data or model is available, or
when deployability to zero-install environments matters more than maximum
sensitivity.

\subsection{Applications}
\label{sec:apps}

The natural applications fall into two classes.

\paragraph{First-stage filters in classical pipelines.} Any search tool
that currently performs $\Omega(N)$ candidate generation (BLAST, DIAMOND,
MMseqs2) can be fronted with the shader kernel to reduce the candidate
set by several orders of magnitude before committing to alignment. This
trades a modest recall penalty (missed at most $5$--$7\%$ of true homologs
at twilight-zone divergence) for a potentially $10^3$--$10^4\times$
reduction in candidates inspected.

\paragraph{Client-side and edge deployments.} The full pipeline can run in
a web browser via WebGL~2, or on a phone via OpenGL~ES; no server-side
alignment service is required. This is relevant to clinical settings
where patient sequence data cannot be uploaded to remote servers, to field
sequencing on nanopore devices with limited connectivity, and to
educational tools where installing a full bioinformatics environment is
not an option.

\subsection{Future work}

Natural extensions, in rough order of effort:
\begin{itemize}
    \item Replacing the global FFT with a sliding-window (``genomic
    spectrogram'') that restores shift-invariance and tolerance to indels.
    \item Evaluation on standard benchmarks (SCOP superfamily retrieval,
    Pfam family assignments, UniRef50 recall) and direct comparison with
    BLAST and DIAMOND on real data.
    \item Actual GLSL implementation and measurement of throughput on
    consumer and mobile GPUs, including a WebGL~2 demonstration.
    \item A learned refinement of the embedding (supervised fine-tuning of
    the per-channel weighting) to close the remaining recall gap against
    exhaustive Jaccard on real data.
    \item Integration with an approximate-nearest-neighbour index
    \citep{malkov2020,johnson2021} when the embedding alone cannot index
    a database of $10^9$+ entries on a single device.
\end{itemize}

%==============================================================================
\section{Conclusion}
\label{sec:conclusion}
%==============================================================================

We described a homology-search architecture that maps sequences to
fixed-dimension spectral coordinates and evaluates similarity as a single
fragment-shader pass. The embedding preserves homolog--non-homolog
separability (ROC-AUC $\geq 0.97$) through the classical twilight zone for
both DNA and protein alphabets. A CPU reference implementation of the
shader kernel runs at $1.1$~ms per query against a $10^5$-sequence
database---three to four orders of magnitude faster than exhaustive $k$-mer
Jaccard and five to seven orders faster than extrapolated exhaustive
Smith--Waterman at the same scale. When combined with a Smith--Waterman
rerank of the top twenty candidates, the pipeline recovers $95\%$ of true
family members on a $1000$-sequence protein benchmark. The method is
positioned as a first-stage filter in front of conventional alignment
tools, not as a replacement; its principal advantages are simplicity of
arithmetic, minimal memory footprint and portability to any GPU or
browser. All code, raw numerical results and figure-generation scripts
are released with the paper so that every number can be reproduced and
extended.

%==============================================================================
\paragraph{Data and code availability.} The JSON result files and the
Python experiment scripts that generated them are in the
\texttt{results/}\ and \texttt{experiments/}\ subdirectories accompanying
this paper. Figures are rebuilt by
\texttt{experiments/make\_figures.py}. No proprietary or non-public
datasets are used.

\paragraph{Author contributions.} Conception, design, implementation,
experiments, analysis, and writing: K.F.S.

\paragraph{Conflicts of interest.} None declared.

%==============================================================================

\bibliography{references}

\end{document}